% §5.3 The landscape at a glance --- blueprint chapter 5, part 05-table.tex (Q-0288; moved forward
% to right after 5.2 in Q-0290, review flag S4: the table, with forward references, and the statement
% of thm:wedge-axis are the section's destination; the theorem's proof is 5.8, 05-wedge-axis.tex,
% after every node it reads). Included from 05-landscape.tex; the conventions are in that file's
% header. The table is built from the nodes, each cell citing the statements it summarizes; every
% cell was compared with its statements for lost hypotheses (WRITING.md § 5, "a summary drops
% hypotheses"). The table is for d >= 2, the dimension in which the general linear Loewner cell and
% the stable-order cells are stated. The one node is thm:wedge-axis; the remark on positivity is the
% plan's (notes/PLAN-article-structure.md, "§5"), as far as it is proved, and states no signed family
% (the author, 2026-10-04).
% Q-0309 (fix group 2, review flag 10): the formal statement of thm:wedge-axis moved to 5.8
% (05-wedge-axis.tex), where its vocabulary exists, before its proof; this subsection keeps the table
% and an informal statement of the theorem (the paragraph "The general linear row along the chain").

\subsection{The landscape at a glance}\label{sec:landscape-table}

Table~\ref{tab:landscape} answers, for each choice of covariance group and of the pairs of apertures
along which the record must be a further smoothing, which kernels are available. It has the groups of \S\ref{sec:scale-shape}
as rows and three orders as columns, for $d \ge 2$, and every cell has an entry, citing the statements
it summarizes in \S\S\ref{sec:invariant-wedges}--\ref{sec:smaller-groups} or in the appendices. The
cells use these terms.
\begin{itemize}
\item The \emph{ray order} relates $\Sigma$ to its multiples $c\,\Sigma$, $c \ge 1$, and $0$ to
everything.
\item The \emph{stable order} $\le_\alpha$, $0 < \alpha < 2$, relates two definite matrices when the
relative step $M = \Sigma^{-1/2}\Sigma'\Sigma^{-1/2}$ has $\lambda_{\min}(M)^{d+\alpha} \ge \det M$, and
adds the origin and the rays (Proposition~\ref{prop:similarity-stable-order-families}(1)). It lies
between the rays and the Loewner order.
% Round 3 (content review), C section 7 (VERIFY-round3-landscape-orbits-paper.md item 30): the pointer
% for "step set", which this page uses before Proposition 4.2.
It is the order of the step set (Proposition~\ref{prop:invariant-orders}, \S\ref{sec:affine-stable-cone})
of the stable orbit of index $\alpha$, the linear images $(\Sigma^{1/2})_*S_\alpha(I)$ of the isotropic $\alpha$-stable law, whose
transform is $e^{-|\omega|^\alpha}$ (Lemma~\ref{lem:elliptical-stable} and
Proposition~\ref{prop:stable-step-cone}). At $\alpha = 1$ that law is the kernel of the Poisson scale
space at unit scale (\S\ref{sec:classical-prior-art}). And it is the order generated by the cone of directions
$C_\alpha := \{X \in \PSD(d) : (d+\alpha)\,\lambda_{\min}(X) \ge \tr X\}$
(Proposition~\ref{prop:stable-order-generated}); $C_0$ gives the rays and $C_\infty := \PSD(d)$ the
Loewner order.
\item The \emph{frame} is the coordinate frame, the axes of the sampling grid along which the
positive diagonal group $\mathrm D_+(d)$ rescales. The \emph{frame's shapes} are the positive definite
diagonal matrices, the apertures whose axes are those of the frame, and they are the orbit of $I$
under $\mathrm D_+(d)$.
% Round 2 (presentation review), item 4 (VERIFY-round2-structure-paper.md § 4.1): D_s and
% Delta(Sigma), defined as in Propositions F.2 and F.4; Corr(d) is not printed in the table.
For $s \in (0,\infty)^d$ write $D_s := \diag(s_1,\dots,s_d)$, so that the frame's shapes are the
$D_s^2$, and write $\Delta(\Sigma) := \diag(\Sigma_{11},\dots,\Sigma_{dd})$ for the diagonal of $\Sigma$.
% Round 2 (presentation review), item 4: the term B-self-decomposable and sym B, as in
% Proposition 4.10.
\item The \emph{exponent cone} $\mathfrak k_\alpha$ of the stable law is the set of matrices $B$ for which
$S_\alpha(I)$ is its own image under $e^{-tB}$ convolved with a residual, for every $t > 0$. A law with
this property is \emph{$B$-self-decomposable}, or operator-self-decomposable with exponent $B$
(Proposition~\ref{prop:stable-exponent-cone}). The symmetric part of $B$ is
$\mathrm{sym}\,B := \tfrac12(B + B^{\tp})$.
\item \emph{Anisotropy} is measured by the exponent spread $\operatorname{es}$ of the non-Gaussian part,
the least $K \ge 1$ with $\psi(Q\omega) \le K\,\psi(\omega)$ for every $Q \in \Orth(d)$. It is set
against the eccentricity $\operatorname{ecc}(\Sigma)$ of the aperture, the ratio of its extreme
eigenvalues. Both are defined again where they are used, in Proposition~\ref{prop:similarity-stable-order-kernels}.
\item A non-Gaussian part is \emph{shape-blind} when it is rotation invariant at every $\Sigma$, so that
it sees the size of the aperture and not its shape. It is shape-blind \emph{at the origin} when its
\Levy{} measure has the same rate near the origin in every direction.
\item The profile entries concern a rotation-invariant law with no Gaussian part whose \Levy{}
density at $x$ is $k(|x|)\,|x|^{-d}$, with $k$ nonincreasing, its \emph{profile}, and the law is
written $\nu_k$. They use the
dilation ratio $R(\lambda) := \inf\{k(r/\lambda)/k(r) : r > 0,\ k(r) > 0\}$, $\lambda \ge 1$, and the
lower exponent $\beta_-(k) := \inf_{\lambda > 1}\log R(\lambda)/\log\lambda$
(Proposition~\ref{prop:orbit-step-cone} and Lemma~\ref{lem:dilation-ratio}).
% Round 2 (presentation review), item 4 (§ 4.2): the gloss of R(lambda) >= lambda^alpha. With
% s = r/lambda it reads k(s) s^alpha >= k(r) r^alpha for s <= r.
The condition $R(\lambda) \ge \lambda^\alpha$ for every $\lambda \ge 1$ says that $r^\alpha k(r)$ is
nonincreasing: towards the origin the profile grows at least like $r^{-\alpha}$.
\end{itemize}
The one-parameter row lies outside the scheme, a curve not being an order on $\PSD(d)$, and is listed
below the others.

\begin{table}[tp]
\centering
\small
\setlength{\tabcolsep}{4pt}
\renewcommand{\arraystretch}{1.15}
\begin{tabular}{@{}p{0.12\textwidth}p{0.27\textwidth}p{0.28\textwidth}p{0.27\textwidth}@{}}
\toprule
$G$ & rays & stable order $\le_\alpha$, $0 < \alpha < 2$ & Loewner order \\
\midrule
$\GL(d)$ &
the orbits $(\Sigma^{1/2})_*\mu_I$ of the rotation-invariant self-decomposable laws $\mu_I \ne \delta_0$,
one family for each law (Propositions~\ref{prop:ray-families} and~\ref{prop:gw-ray-families}(3)) &
% Round 3 (content review), E i and C-R3 (the author's decision D10): the cell is complete,
% Proposition 4.22(2) with C5.
the orbits of exactly the laws $N(0, aI)$, $a > 0$, and $N(0, aI) * \nu_k$, $a \ge 0$, with $k$
right-continuous and $\beta_-(k) \ge \alpha$; among the stable orbits, exactly those of index $\gamma \ge \alpha$
(Theorem~\ref{thm:wedge-axis}(3), Proposition~\ref{prop:wedge-axis-profiles},
Lemma~\ref{lem:dilation-ratio}(3)) &
the affine Gaussian scale spaces only (Theorem~\ref{thm:affine-rigidity}) \\
\addlinespace
similarities $\mathrm S(d)$ &
continuous covariant fields of symmetric self-decomposable laws over the shapes, with orbits indexed by
the normalized spectrum; the isotropic members are the rotation-invariant self-decomposable laws
$\ne \delta_0$, and none fixes its family (Proposition~\ref{prop:similarity-ray-families}) &
every family's kernels are symmetric infinitely divisible with
$\operatorname{es}(\psi_\Sigma) \le 4\operatorname{ecc}(\Sigma)^{(d+\alpha)/\alpha}$; contains the restrictions of
the $\GL(d)$ cell and of the Loewner cell, and a family in neither
(Proposition~\ref{prop:similarity-stable-order-kernels}) &
% Round 3 (content review), A.R9: F.3(4) and its witness F.3(3)(ii).
the isotropic members of the ray cell, each with a shape-blind family
(Proposition~\ref{prop:similarity-ray-families}(3)--(4)); every family's kernels are infinitely divisible with
$\operatorname{es}(\psi_\Sigma) \le 4\operatorname{ecc}(\Sigma)$
(Proposition~\ref{prop:similarity-stable-order-kernels}(2)); a uniform bound when the \Levy{} measures
at the thinnest apertures stay comparable to the isotropic member's, and none in general
(Propositions~\ref{prop:similarity-general-bounded-shape}, \ref{prop:similarity-cascade-families}
and~\ref{prop:similarity-degenerate-family}); for a profile bounded at the origin, shape-blind at the
origin (Proposition~\ref{prop:similarity-general-index-zero}) and need not be beyond it
% Round 2 (presentation review), item 4 (§ 4.3): "not beyond it" narrowed to "need not be beyond
% it"; one example family backs it, and the shape-blind families of the cell are rotation invariant.
(Proposition~\ref{prop:similarity-line-family}(3)); not every family is a Gaussian convolved with a
shape-blind part (Proposition~\ref{prop:similarity-stable-family}(3)) \\
\addlinespace
positive diagonal $\mathrm D_+(d)$ &
on the frame's shapes the images $(D_s)_*\mu_I$ of any self-decomposable $\mu_I \ne \delta_0$, not
necessarily products; orbits in $\PD(d)$ indexed by the correlation matrices; a rotation-invariant member fixes
no family off the frame's shapes (Proposition~\ref{prop:diagonal-ray-families}) &
on the frame's shapes $(D_s)_*\mu_I$ with $\mu_I$ operator-self-decomposable for every diagonal exponent
in $\mathfrak k_\alpha$, a necessary condition, and whether every such law is the isotropic member of a
family is open; strictly between the restricted Loewner cell and the ray cell, and not the union of the
restrictions of the Loewner and $\GL(d)$ cells; off the frame's shapes not described
(Proposition~\ref{prop:diagonal-stable-order-families}) &
on the frame's shapes the products $(s_1)_*\nu_1 \otimes \dots \otimes (s_d)_*\nu_d$ of self-decomposable
laws on the line (Propositions~\ref{prop:fixed-frame-families} and~\ref{prop:diagonal-ray-families}(3));
off the frame's shapes one extension, $(\Delta(\Sigma)^{1/2})_*(\nu_1 \otimes \dots \otimes \nu_d)$, the
families not described \\
\midrule
% Round 3 (content review), C-S6: what "with rotations" meant.
\multicolumn{4}{@{}p{0.98\textwidth}@{}}{\emph{One linear flow} $\{e^{sB}\}$, $B$ with a skew part allowed, along the
curve $e^{sB}\Sigma_0e^{sB^{\tp}}$: the families are the $B$-self-decomposable laws
(Proposition~\ref{prop:one-parameter-families}); for a contracting flow a rotation-invariant non-Gaussian
member exists exactly when $\tr B < (d+2)\lambda_{\min}(\mathrm{sym}\,B)$
(Proposition~\ref{prop:anisotropy-budget}).} \\
\bottomrule
\end{tabular}
\caption{The affine landscape for $d \ge 2$: the families of each group along each order. The four
corners of the scheme (Definition~\ref{def:gw-cascade-family}) are the cells of the first two rows in
the ray and the Loewner columns. Each cell cites the statements it summarizes. The cells of the general
linear group are complete, and \S\ref{sec:conclusion} says how far the others go.}
% Round 3 (content review), A.R1 and C-R2 (PLAN-review-response-paper.md R3.1, "taken unless vetoed"):
% the complete cells named in the caption, three since D10.
\label{tab:landscape}
\end{table}

\paragraph{The inclusions.} A family of a larger group or a larger order restricts to every smaller
datum, with the same kernels from scale zero (Lemma~\ref{lem:gw-kernels-from-origin}(1)), so in each column
the cells grow from the general linear row to each smaller group's row, and along each row towards the
rays. These inclusions are known to be strict:
\begin{itemize}
\item along the general linear row, at every step of the chain up to $\alpha = 2$
(Theorem~\ref{thm:wedge-axis}(3));
\item from the $\GL(d)$ cell and the restricted Loewner cell to the cell (similarities, stable order)
(Proposition~\ref{prop:similarity-stable-order-kernels}(3));
\item from the restricted Loewner cell to the stable cell, and from there to the ray cell, in the diagonal
row (Proposition~\ref{prop:diagonal-stable-order-families}(4)), and so from the restricted (diagonal,
Loewner) cell to the (diagonal, rays) cell (Proposition~\ref{prop:diagonal-ray-families}(4)).
\end{itemize}
% Round-1 external review (Gemini; PLAN-review-response-paper.md, D3): the paragraphs "The isotropic
% members" (moved to section 1, "Smaller groups") and "The general linear row along the chain" (said by
% section 1 and section 4.7) are folded into section 1.

% Q-0317 (the blind register review, second half flag 25): rem:positivity moved to section 3.4,
% beside the rigidity it is about.
